\chapter{Contrats perpétuels}\label{ch:m3:perpetual-futures}

À minuit, à huit heures et à seize heures, en UTC, les détenteurs de \omterm{def:m3:perpetual-futures:perp}{contrats perpétuels} sur crypto-actifs se versent
de l'argent. Si le contrat s'est échangé au-dessus du prix au comptant au cours des huit
heures précédentes, les acheteurs paient les vendeurs~; en dessous, ce sont les vendeurs qui paient les acheteurs. Le contrat n'expire jamais, ne
donne jamais lieu à livraison et n'a jamais besoin d'être roulé, et pourtant il reste à une fraction de pourcent du prix au comptant, car
ce paiement périodique rend coûteux le fait de détenir le côté qui pousse le prix à s'écarter. Le \omterm{def:m3:perpetual-futures:perp}{contrat perpétuel}
est l'instrument de levier principal du marché des crypto-actifs~: parmi l'ensemble des contrats à terme sur bitcoin cotés et négociés au premier semestre 2023,
les perpétuels représentaient 75~\% du volume quotidien moyen de 27~milliards de dollars et 94~\% de
l'encours ouvert moyen de 8~milliards de dollars, selon le décompte d'Ackerer, Hugonnier et Jermann. Ce chapitre le détaille~:
l'indice qu'il suit et le \omterm{def:m3:perpetual-futures:index}{prix de marquage} auquel il est valorisé, la prime et la formule qui la
transforme en paiement de funding, les trois manières de libeller le contrat et l'opération qui encaisse le
funding.

\section{Un contrat à terme qui n'expire jamais}

Un contrat à échéance (livre~1, chapitre~18) converge vers le prix au comptant parce qu'à l'échéance il devient le prix au comptant~: la
position longue reçoit le sous-jacent, ou sa valeur en espèces, au prix de règlement-livraison final. Entre maintenant et
l'échéance, le contrat peut se négocier avec une base par rapport au prix au comptant, que l'opération de \omterm{def:m3:forward-curves-storage-and-convenience-yield:carry}{cash-and-carry} maintient près du
coût de portage. Supprimer l'échéance, et rien ne force la convergence. Le \omterm{def:m3:perpetual-futures:perp}{contrat perpétuel} remplace la livraison par un flux de
paiements.

\begin{definition}[Contrat perpétuel, taux de financement, intervalle de financement]\label{def:m3:perpetual-futures:perp}
Un \emph{contrat perpétuel}\index{contrat perpétuel} (perpetual future) est un contrat à terme sans échéance, réévalué en continu,
dont les détenteurs échangent des paiements périodiques qui poussent son prix vers un indice du
prix au comptant du sous-jacent. Le paiement par unité de valeur de position est le \emph{taux de
financement}\index{taux de financement} (funding rate), positif lorsque les acheteurs paient les vendeurs~; la période entre deux
paiements est l'\emph{intervalle de financement}\index{intervalle de financement} (funding interval), généralement de huit heures.
\end{definition}

Dans ce chapitre, $f$ désigne le \omterm{def:m3:perpetual-futures:perp}{taux de financement} par intervalle, et une position de $q$ unités du
sous-jacent détenue à un instant de financement verse $q\,M f$, où $M$ est le \omterm{def:m3:perpetual-futures:index}{prix de marquage} défini plus bas~: une position longue paie
lorsque $f > 0$ et reçoit lorsque $f < 0$. Si le perpétuel se négocie au-dessus du prix au comptant, les positions longues paient, ce qui
incite les arbitragistes à vendre le perpétuel et à acheter le comptant, en encaissant le funding tout en étant couverts~; leurs
ventes font redescendre le perpétuel. Le funding joue pour le perpétuel le rôle que la convergence joue pour un contrat à échéance.

\begin{proposition}[Le perpétuel comme contrat à terme roulé en continu]\label{prop:m3:perpetual-futures:roll}
Considérons une position acheteuse sur le perpétuel et une position vendeuse au comptant de même taille, et supposons que le prix du perpétuel soit égal
à l'indice au début et à la fin. Le P\&L total de la position est l'opposé de la somme des paiements de funding. Un contrat à
échéance acheté à une base $b$ par rapport au comptant et conservé jusqu'à son échéance rapporte $-b$ par rapport au comptant. Le funding est
donc la base du perpétuel payée par tranches~: un perpétuel dont le funding moyen vaut $\bar f$ par intervalle
se comporte comme un contrat à terme roulé à une base de $\bar f$ par intervalle.
\end{proposition}

\begin{proof}
Les prix étant égaux aux deux extrémités, les jambes de prix se compensent et seuls subsistent les flux de funding. Pour le
contrat à échéance, la convergence fait passer son prix de $S(1+b)$ à $S_T$ alors que le comptant passe de $S$ à $S_T$, soit un écart de $-bS$.
\end{proof}

\begin{dated}{2026-09}{Historique d'un perpétuel et son funding par défaut}\label{dat:m3:perpetual-futures:bitmex}
Les données sur les instruments de BitMEX donnent le 13 mai 2016 comme date d'admission à la cote de son swap perpétuel XBTUSD, un \omterm{def:m3:perpetual-futures:types}{contrat
inverse} réglé en bitcoin avec un \omterm{def:m3:perpetual-futures:perp}{intervalle de financement} de huit heures~; le 24 septembre 2026, son \omterm{def:m3:perpetual-futures:perp}{taux de financement}
était de 0.01~\% par intervalle. Binance verse le funding à 00:00, 08:00 et 16:00 UTC par défaut, avec une composante d'intérêt
de 0.01~\% par intervalle (0.03~\% par jour) pour la plupart des contrats.
\end{dated}

\section{Prix de l'indice et prix de marquage}

Un perpétuel a besoin de deux prix de référence, en plus de son dernier prix échangé.

\begin{definition}[Prix de l'indice, prix de marquage]\label{def:m3:perpetual-futures:index}
Le \emph{prix de l'indice}\index{prix de l'indice} (index price) d'un perpétuel est une moyenne pondérée du prix au comptant du
sous-jacent sur plusieurs plateformes constitutives, selon des règles qui limitent l'influence de chacune d'elles. Le \emph{prix de
marquage}\index{prix de marquage} (mark price) est le prix auquel la plateforme valorise les positions ouvertes pour le P\&L latent, la marge
et la liquidation~; il est construit à partir de l'indice et d'une base lissée, et non à partir de la dernière transaction.
\end{definition}

L'indice est la cible vers laquelle le funding ramène le contrat~; le \omterm{def:m3:perpetual-futures:index}{prix de marquage} détermine qui est liquidé. Les deux sont conçus
contre la manipulation. Un trader capable de faire bouger le cours sur une seule plateforme au comptant, ou le dernier prix
échangé du perpétuel lui-même, pendant quelques secondes, pourrait sinon déclencher des liquidations et en tirer profit
(\cref{ch:m3:margin-liquidation-and-loss-allocation}). Un indice qui plafonne chaque composante près de la médiane
de l'ensemble et écarte celles qui cessent de se mettre à jour coûte plus cher à déplacer~; un \omterm{def:m3:perpetual-futures:index}{prix de marquage} fondé sur une
médiane de plusieurs estimations ignore une transaction isolée et aberrante.

\begin{dated}{2026-09}{Règles de l'indice et du prix de marquage sur une grande plateforme}\label{dat:m3:perpetual-futures:mark}
Les contrats à terme à marge en USD de Binance calculent l'indice comme une somme pondérée de prix au comptant
sur les plateformes constitutives, qu'elle cite comme incluant Binance, OKX, Coinbase, Kraken, Bybit et d'autres,
ainsi que des plateformes décentralisées. Une composante dont le dernier cours s'écarte de plus de 3~\% de la médiane de
l'ensemble des sources est plafonnée à 1.03 ou 0.97 fois la médiane (1~\% pour certaines paires comme BTCUSDT)~; une composante qui
ne répond pas ou n'a pas été mise à jour depuis cinq minutes reçoit un poids nul. Le \omterm{def:m3:perpetual-futures:index}{prix de marquage} est la médiane de trois
prix~: l'indice multiplié par $1 + f \times (\text{temps jusqu'au prochain financement}/\text{période de financement})$ avec le dernier
\omterm{def:m3:perpetual-futures:perp}{taux de financement} $f$~; l'indice plus une moyenne mobile sur 30~secondes de la base~; et le dernier cours du contrat.
\end{dated}

\section{Indice de prime et formule de funding}

\begin{definition}[Indice de prime]\label{def:m3:perpetual-futures:premium}
L'\emph{indice de prime}\index{indice de prime} (premium index) mesure l'écart entre le carnet d'ordres du perpétuel et l'indice~:
où $I$ est l'indice et $B^{\mathrm{imp}}$, $A^{\mathrm{imp}}$ les prix moyens auxquels un notionnel donné
(le notionnel d'impact) serait vendu sur le carnet côté acheteur ou acheté sur le carnet côté vendeur,
\[
p = \frac{\max(0,\, B^{\mathrm{imp}} - I) - \max(0,\, I - A^{\mathrm{imp}})}{I}.
\]
\end{definition}

La prime est nulle tant que l'indice se situe entre le cours acheteur et le cours vendeur d'impact~; elle est positive lorsque même les
cours acheteurs sont au-dessus de lui, négative lorsque même les cours vendeurs sont en dessous. Utiliser les prix d'impact plutôt que les meilleurs cours rend
coûteux de le déplacer avec de petits ordres. La plateforme l'échantillonne tout au long de l'intervalle et moyenne les échantillons avec des poids
croissants dans le temps, de sorte que la prime proche de l'instant de financement compte le plus~: avec $n$ échantillons $p_1,\dots,p_n$,
$\bar p = \sum_k k\,p_k \big/ \sum_k k$.

\begin{proposition}[La règle de funding et sa zone morte]\label{prop:m3:perpetual-futures:rule}
Soient $i$ la composante d'intérêt par intervalle et $c$ la borne de troncature. Le \omterm{def:m3:perpetual-futures:perp}{taux de financement}
\[
f = \bar p + \operatorname{clamp}(i - \bar p,\, -c,\, c)
\]
vaut $i$ dès que $|\bar p - i| \le c$, vaut $\bar p - c$ lorsque $\bar p > i + c$ et $\bar p + c$ lorsque
$\bar p < i - c$. Avec $i = 0.01\%$ et $c = 0.05\%$, le funding vaut exactement $0.01\%$ pour toute prime moyenne
comprise entre $-0.04\%$ et $+0.06\%$.
\end{proposition}

\begin{proof}
Si $-c \le i - \bar p \le c$, la troncature renvoie $i - \bar p$ et $f = i$. Si $i - \bar p < -c$, elle renvoie $-c$~;
si $i - \bar p > c$, elle renvoie $c$.
\end{proof}

\begin{omfigure}
\begin{tikzpicture}
\begin{axis}[width=11cm,height=5cm,
  xlabel={indice de prime moyenne, points de base par intervalle},ylabel={taux de financement, pb},
  tick label style={font=\small},label style={font=\small},
  xmin=-30,xmax=30,ymin=-30,ymax=30,grid=major,grid style={gray!25},
  legend columns=2,legend style={font=\footnotesize,at={(0.5,-0.3)},anchor=north,draw=none,fill=none,/tikz/every even column/.append style={column sep=8pt}},
  table/col sep=comma]
\addplot[omThm,very thick] table[x=premium_bp,y=funding_bp]{figdata/markets-3/17-perpetual-futures/rule.csv};\addlegendentry{taux de financement}
\addplot[gray,dashed,domain=-30:30,samples=2] {x};\addlegendentry{financement égal à la prime}
\draw[omDef,thick,{Stealth[length=2mm]}-{Stealth[length=2mm]}] (axis cs:-4,6) -- (axis cs:6,6);
\node[font=\footnotesize,anchor=south,text=omDef] at (axis cs:1,6.5) {zone morte};
\end{axis}
\end{tikzpicture}

\omcaption{La règle de financement de la \cref{prop:m3:perpetual-futures:rule} avec une composante d'intérêt de 1 point de base
et une troncature de 5 points de base par intervalle~: le financement reste à 1 point de base pour toute prime moyenne comprise entre $-4$ et
$+6$ points de base et suit la prime, décalé de 5 points de base vers la composante d'intérêt, en dehors de cette zone.
Données~: le tutoriel du chapitre.}\label{fig:m3:perpetual-futures:rule}
\end{omfigure}

La zone morte a une conséquence qui surprend les débutants~: sur les marchés calmes, le \omterm{def:m3:perpetual-futures:perp}{taux de financement} n'est pas la prime, mais
la composante d'intérêt, 0.01~\% toutes les huit heures, soit $0.01\% \times 3 \times 365 = 10.95\%$ par an, versés par les positions longues aux positions courtes.
La composante d'intérêt est censée traduire l'écart entre les taux d'intérêt de la devise de cotation et ceux du sous-jacent,
comme pour les points de terme du livre~2~; en pratique, c'est un paramètre fixe,
et, pendant l'essentiel d'une année calme, une position courte sur \omterm{def:m3:perpetual-futures:perp}{contrat perpétuel}, couverte au comptant, l'encaisse.

\begin{dated}{2026-09}{La formule de financement sur une grande plateforme}\label{dat:m3:perpetual-futures:formula}
Le \omterm{def:m3:perpetual-futures:perp}{taux de financement} de Binance vaut $[\bar p + \operatorname{clamp}(i - \bar p, 0.05\%, -0.05\%)]/(8/N)$ pour un
intervalle de $N$ heures, avec $i = 0.01\%$ par intervalle pour la plupart des contrats. L'\omterm{def:m3:perpetual-futures:premium}{indice de prime} utilise les cours acheteur et vendeur d'impact calculés pour
un notionnel de marge d'impact de 200 USDT divisé par le taux de marge initiale au levier maximal, échantillonnés toutes les 5 secondes
et moyennés avec des poids linéairement croissants (5\,760 échantillons sur huit heures). Le financement est encadré par un plafond et un plancher égaux à
$\pm 0.75$ fois le ratio de marge de maintenance pour les contrats majeurs, tels que BTCUSDT, et à $\pm 2\%$ pour les autres.
\end{dated}

\begin{omfigure}
\begin{tikzpicture}
\begin{axis}[width=11cm,height=5cm,
  xlabel={jour},ylabel={points de base par intervalle},
  tick label style={font=\small},label style={font=\small},
  xmin=0,xmax=30,ymin=-10,ymax=16,grid=major,grid style={gray!25},
  legend columns=2,legend style={font=\footnotesize,at={(0.5,-0.3)},anchor=north,draw=none,fill=none,/tikz/every even column/.append style={column sep=8pt}},
  table/col sep=comma]
\addplot[gray,thick,mark=*,mark size=0.8pt,only marks] table[x=day,y=premium_bp]{figdata/markets-3/17-perpetual-futures/month.csv};\addlegendentry{prime moyenne}
\addplot[omThm,very thick,const plot] table[x=day,y=funding_bp]{figdata/markets-3/17-perpetual-futures/month.csv};\addlegendentry{taux de financement}
\end{axis}
\end{tikzpicture}

\omcaption{Trente jours synthétiques, par intervalles de huit heures~: calme, une hausse pendant laquelle le \omterm{def:m3:perpetual-futures:perp}{contrat perpétuel} se traite cher, une
forte baisse pendant laquelle il se traite à bon marché, puis à nouveau le calme. Pendant les périodes calmes, la prime moyenne oscille dans la zone morte et
le financement reste fixé à 1 point de base~; lors de la hausse, les acheteurs paient jusqu'à environ 10 points de base par intervalle~; lors de la baisse,
les vendeurs paient. Échantillons de prime simulés toutes les 5 secondes~; la règle de la plateforme appliquée à l'identique. Données~: le
tutoriel du chapitre.}\label{fig:m3:perpetual-futures:month}
\end{omfigure}

\section{Contrats linéaires, inverses et quanto}

Un \omterm{def:m3:perpetual-futures:perp}{contrat perpétuel} sur bitcoin peut être libellé de trois façons, et ce choix modifie l'exposition du détenteur.

\begin{definition}[Contrats linéaires, inverses et quanto]\label{def:m3:perpetual-futures:types}
Un \emph{contrat linéaire}\index{contrat linéaire} est coté et réglé dans la devise de cotation (un dollar ou un
\omterm{def:m3:blockchains-for-traders:stable}{stablecoin}), son P\&L est linéaire dans le prix~: $q\,(S - E)$ pour $q$ unités achetées à $E$. Un \emph{contrat
inverse}\index{contrat inverse} a un notionnel fixe exprimé dans la devise de cotation, mais il est margé et réglé dans la
cryptomonnaie sous-jacente, de sorte que $N$ contrats d'un dollar chacun achetés à $E$ rapportent $N\,(1/E - 1/S)$ unités de cryptomonnaie. Un
\emph{contrat quanto}\index{contrat quanto} est coté sur un actif mais margé et réglé dans un autre à un taux fixe
par unité de variation du prix, de sorte que son P\&L dans l'actif de règlement vaut $N\,m\,(S - E)$ pour un multiplicateur $m$,
quel que soit le prix de l'actif de règlement.
\end{definition}

\begin{dated}{2026-09}{Conditions des contrats chez BitMEX}\label{dat:m3:perpetual-futures:terms}
Les données d'instruments de BitMEX au 24 septembre 2026 décrivent XBTUSD comme un perpétuel inverse, un dollar US de bitcoin par
contrat, réglé en bitcoin~; ETHUSD, admis à la cote le 1\up{er} août 2018, comme un perpétuel quanto coté en dollars et réglé en
bitcoin, avec un multiplicateur de 100 satoshis (0.000001 bitcoin) par contrat et par dollar de prix de l'ether~; et XBTUSDT,
admis à la cote en novembre 2021, comme un perpétuel linéaire réglé en USDT. Les trois avaient un \omterm{def:m3:perpetual-futures:perp}{intervalle de financement} de huit heures.
\end{dated}

\begin{proposition}[Le P\&L en cryptomonnaie d'un contrat inverse]\label{prop:m3:perpetual-futures:inverse}
Le P\&L en cryptomonnaie d'une position longue inverse de $N$ contrats d'un dollar, $N(1/E - 1/S)$, est concave en $S$~: sa
dérivée $N/S^2$ diminue à mesure que le prix monte. Une position longue gagne au plus $N/E$ unités de cryptomonnaie, quelle que soit la hausse du prix, et perd
sans limite lorsque $S \to 0$~; le P\&L de la position courte en est l'image symétrique. En dollars, $N(1/E - 1/S)\,S = N(S/E - 1)$ est
linéaire.
\end{proposition}

\begin{proof}
On dérive deux fois~: $N/S^2 > 0$ et $-2N/S^3 < 0$. Les limites découlent de $1/S \to 0$ et $1/S \to \infty$.
\end{proof}

\begin{omfigure}
\begin{tikzpicture}
\begin{axis}[width=11cm,height=5cm,
  xlabel={prix du bitcoin, USD},ylabel={P\&L, BTC},
  tick label style={font=\small},label style={font=\small},
  xmin=30000,xmax=120000,ymin=-1.1,ymax=1.1,grid=major,grid style={gray!25},
  xtick={30000,60000,90000,120000},xticklabels={30\,000,60\,000,90\,000,120\,000},scaled x ticks=false,
  legend columns=2,legend style={font=\footnotesize,at={(0.5,-0.3)},anchor=north,draw=none,fill=none,/tikz/every even column/.append style={column sep=8pt}},
  table/col sep=comma]
\addplot[omThm,very thick] table[x=price,y=pnl_btc]{figdata/markets-3/17-perpetual-futures/inverse.csv};\addlegendentry{position longue de 60\,000 contrats inverses}
\addplot[omDef,thick,dashed] table[x=price,y=tangent_btc]{figdata/markets-3/17-perpetual-futures/inverse.csv};\addlegendentry{tangente à l'entrée}
\end{axis}
\end{tikzpicture}

\omcaption{P\&L en bitcoin d'une position longue de 60\,000 \omterm{def:m3:perpetual-futures:types}{contrats inverses} d'un dollar, établie à USD~60\,000
(\cref{prop:m3:perpetual-futures:inverse})~: une baisse à 30\,000 coûte un bitcoin entier, une hausse à 120\,000 ne rapporte que
la moitié d'un bitcoin. La courbe se situe sous sa tangente~: une position longue inverse est vendeuse de convexité en cryptomonnaie. Données~: le
tutoriel du chapitre.}\label{fig:m3:perpetual-futures:inverse}
\end{omfigure}

Le \omterm{def:m3:perpetual-futures:types}{contrat inverse} convient à un détenteur qui raisonne en coins~: un mineur ou un fonds dont les comptes sont tenus en bitcoins peut couvrir sa valeur en dollars
en vendant des \omterm{def:m3:perpetual-futures:types}{contrats inverses}, puisqu'une position courte inverse, augmentée des coins déposés en marge, a une valeur en dollars constante.
Un trader qui raisonne en dollars et prend une position longue avec une marge en coins est doublement long~: la position
et le collatéral baissent ensemble. Le \omterm{def:m3:perpetual-futures:types}{contrat quanto} donne une exposition à l'ether réglé en bitcoins, à un nombre fixe
de satoshis par dollar~; son P\&L en dollars est le P\&L en coins multiplié par le cours du bitcoin, de sorte que le détenteur est aussi exposé
au bitcoin et à la corrélation entre les deux, que le financement et le prix du contrat doivent
compenser. Les \omterm{def:m3:perpetual-futures:types}{contrats linéaires} margés en \omterm{def:m3:blockchains-for-traders:stable}{stablecoins} évitent ces deux effets, mais portent en revanche le risque propre du \omterm{def:m3:blockchains-for-traders:stable}{stablecoin}
(\cref{ch:m3:blockchains-for-traders}).

\section{Cash-and-carry et arbitrage de funding}

\begin{definition}[Arbitrage de funding]\label{def:m3:perpetual-futures:funding-arb}
L'\emph{arbitrage de funding}\index{arbitrage de funding} consiste à détenir un \omterm{def:m3:perpetual-futures:perp}{contrat perpétuel} face à une position de sens contraire dans
le sous-jacent (au comptant, ou sur le perpétuel d'une autre plateforme) afin que les variations de prix se compensent et que la position encaisse l'écart de
financement~; il s'agit le plus souvent d'être long au comptant et court sur le perpétuel pour encaisser un financement positif.
\end{definition}

Il s'agit du \omterm{def:m3:forward-curves-storage-and-convenience-yield:carry}{cash-and-carry} du livre~1, chapitre~21, le carry étant versé toutes les huit heures au lieu d'être
réglé à l'échéance. Son rendement dépend de la quantité de financement à encaisser. Schmeling, Schrimpf et Todorov, de la BRI, constatent
que le carry des cryptomonnaies, soit l'écart entre les prix des contrats à terme et ceux du comptant, dépasse en moyenne 10~\% par an, bien plus que le carry
dans les actions, les obligations, les devises ou les matières premières, est extrêmement volatil, atteignant 60~\% par an, et il est au plus haut lorsque les petits
investisseurs recherchent la hausse avec effet de levier et que le capital d'arbitrage est rare~; un carry élevé annonce des krachs ultérieurs.

\begin{dated}{2026-09}{Funding d'un contrat perpétuel majeur, par année}\label{dat:m3:perpetual-futures:years}
D'après l'historique public des \omterm{def:m3:perpetual-futures:perp}{taux de financement} de Binance pour le \omterm{def:m3:perpetual-futures:perp}{contrat perpétuel} BTCUSDT (calculé par l'auteur le 24 septembre 2026), le
\omterm{def:m3:perpetual-futures:perp}{taux de financement} moyen annualisé était de 17.2~\% en 2020, 30.6~\% en 2021, 4.2~\% en 2022, 7.9~\% en 2023, 11.9~\% en 2024, 5.1~\%
en 2025 et 2.9~\% en 2026 à ce jour~; la part des événements de financement négatifs était la plus élevée en 2022 (22~\%) et 2026 (26~\%).
Le plus fort taux unitaire était de 0.3~\% par intervalle en 2020~; depuis 2025, le financement n'a pas dépassé la valeur par défaut de 0.01~\%.
\end{dated}

\begin{omfigure}
\begin{tikzpicture}
\begin{axis}[width=11cm,height=5cm,ybar,bar width=10pt,
  xlabel={année},ylabel={\% par an},
  tick label style={font=\small},label style={font=\small},
  xmin=2019.4,xmax=2026.6,ymin=0,ymax=35,grid=major,grid style={gray!25},
  xtick={2020,2021,2022,2023,2024,2025,2026},xticklabel style={/pgf/number format/1000 sep=},
  legend columns=2,legend style={font=\footnotesize,at={(0.5,-0.3)},anchor=north,draw=none,fill=none,/tikz/every even column/.append style={column sep=8pt}},
  table/col sep=comma]
\addplot[fill=omThm!60,draw=omThm,area legend] table[x=year,y=funding_pct]{figdata/markets-3/17-perpetual-futures/years.csv};\addlegendentry{funding moyen, annualisé}
\addplot[fill=omDef!60,draw=omDef,area legend] table[x=year,y=trade_pct]{figdata/markets-3/17-perpetual-futures/years.csv};\addlegendentry{funding trade, rendement du capital}
\end{axis}
\end{tikzpicture}

\omcaption{Funding annualisé moyen du \omterm{def:m3:perpetual-futures:perp}{contrat perpétuel} BTCUSDT sur Binance, par année civile (2026 jusqu'au 24 septembre), et le
rendement du capital d'une position longue au comptant et courte sur le perpétuel, conservée un an à prix constant, avec 0.30~\% de frais aller-retour et
30.65~\% de marge pour la jambe courte (\cref{prop:m3:perpetual-futures:margin}). Données~: statistiques récapitulatives calculées à partir de
l'historique public des \omterm{def:m3:perpetual-futures:perp}{taux de financement} de Binance~; le trade~: le tutoriel du chapitre.}\label{fig:m3:perpetual-futures:years}
\end{omfigure}

La stratégie présente trois faiblesses. Le financement peut devenir négatif, comme il l'a fait un intervalle sur cinq en 2022, et la position couverte
doit alors payer. Le perpétuel court est margé séparément du comptant, si bien qu'une hausse profitant au comptant
vide la marge du short~; si celle-ci n'est pas reconstituée à temps, le short est liquidé et la position reste
non couverte au pire moment.

\begin{proposition}[Marge pour résister à une hausse]\label{prop:m3:perpetual-futures:margin}
Un perpétuel court de notionnel $Q$ ouvert au prix $E$, avec un taux de marge de maintenance $m$, résiste à une hausse de $x$
(le cours passant à $E(1+x)$) sans liquidation que si sa marge vaut au moins $Q\,[x + m(1 + x)]$~: pour une
hausse de 30~\% et $m = 0.5\%$, soit 30.65~\% du notionnel initial.
\end{proposition}

\begin{proof}
À $E(1+x)$, le short a perdu $Qx$ et doit encore détenir une marge de maintenance égale à $m$ fois son notionnel, soit désormais $Q(1+x)$.
\end{proof}

La troisième faiblesse tient à la plateforme~: le comptant et la marge sont déposés sur des bourses, et le rendement du trade est un spread de crédit
par rapport à ce qu'une bourse défaillante rembourserait (\cref{ch:m3:centralised-exchanges}). Les plateformes qui font compter les avoirs au comptant comme
collatéral pour le short (marge de portefeuille ou marge croisée, \cref{ch:m3:margin-liquidation-and-loss-allocation}) suppriment
l'essentiel du coût de la marge, au prix d'une concentration des deux jambes sur une seule plateforme.

\section{Tutoriel~: funding et contrat inverse}\label{tut:m3:perpetual-futures:funding}

\begin{tutorial}
\textbf{Objectif.} Implémenter l'\omterm{def:m3:perpetual-futures:premium}{indice de prime} et la règle de financement, les appliquer à un mois simulé d'échantillons de prime,
et tracer le P\&L en coins du \omterm{def:m3:perpetual-futures:types}{contrat inverse}. \textbf{État final~:}
\cref{fig:m3:perpetual-futures:rule,fig:m3:perpetual-futures:month,fig:m3:perpetual-futures:inverse} et les chiffres
du problème du week-end.
\begin{tutsteps}
\item \textbf{La formule.} Prime, moyenne pondérée dans le temps et financement, tels que la plateforme les publie.
\omcode{code/firm/perp/firm_perp.py}{42}{56}{L'indice de prime, sa moyenne pondérée dans le temps et le taux de financement.}\label{lst:m3:perpetual-futures:funding}
\item \textbf{Le mois.} \texttt{month\_of\_funding(regimes())} simule 5\,760 échantillons de prime par intervalle.
\item \textbf{Le trade.} La marge nécessaire à un short et le rendement du funding trade.
\omcode{code/markets-3/17-perpetual-futures/python/m3_perp.py}{48}{58}{Marge pour résister à une hausse et rendement du capital du funding trade.}\label{lst:m3:perpetual-futures:trade}
\item \textbf{Exécuter} \texttt{fig\_perp.py} pour obtenir les données des graphiques.
\end{tutsteps}
\textbf{Pour aller plus loin.} Faire payer le financement toutes les heures au lieu de toutes les huit heures et observer comment la zone morte
évolue~; faire supporter au trade le financement de 2022 avec ses intervalles négatifs~; valoriser l'exposition du \omterm{def:m3:perpetual-futures:types}{contrat quanto} à la corrélation en simulant
conjointement le bitcoin et l'ether.
\end{tutorial}

\section{Construction~: le moteur de contrats perpétuels}\label{bld:m3:perpetual-futures:perp}

\begin{build}
\bfield{Rôle} Tout ce que la mini-entreprise fait avec les perpétuels (cotation, couverture, funding trade, moteur
de liquidation du \cref{ch:m3:margin-liquidation-and-loss-allocation}) exige que l'indice, le \omterm{def:m3:perpetual-futures:index}{prix de marquage} et le financement de la plateforme soient
reproduits à l'identique, ainsi que le P\&L dans la bonne devise pour chaque type de contrat.
\bfield{Interface} \texttt{index\_price(sources, cap, stale\_s)}~; \texttt{impact\_price(levels, notional)}~;
\texttt{premium\_index(impact\_bid, impact\_ask, index)}~; \texttt{average\_premium(samples)}~;
\texttt{funding\_rate(avg\_premium, interest, clamp, cap)}~; \texttt{mark\_price(index, last\_funding,
frac\_to\_funding, basis\_ma, last\_trade)}~; \texttt{funding\_payment}~; \texttt{pnl\_linear}, \texttt{pnl\_inverse},
\texttt{pnl\_quanto}.
\bfield{Règles} Paramètres par plateforme et par contrat, lus dans leur documentation et versionnés~; aucune source, aucun indice (une
erreur, jamais une valeur périmée)~; taux par intervalle.
\bfield{Tests d'acceptation} \texttt{code/firm/perp/tests/}~: plafonnement d'une valeur aberrante et exclusion d'une source périmée~;
prix d'impact sur deux niveaux~; prime nulle à l'intérieur de la fourchette d'impact~; financement à l'intérieur et à l'extérieur de la zone morte, et au
plafond~; la médiane du \omterm{def:m3:perpetual-futures:index}{prix de marquage}~; le P\&L des trois types de contrat.
\bfield{Pour aller plus loin} \omterm{def:m3:perpetual-futures:perp}{Intervalles de financement} horaires et dynamiques~; reproduction de l'historique publié par la plateforme à partir
de ses données de carnet d'ordres~; ajustement de la corrélation du \omterm{def:m3:perpetual-futures:types}{contrat quanto}.
\end{build}

\begin{omsources}
\item Binance, ``Introduction to Binance Futures Funding Rates'' et ``How to Calculate the Mark Price and the Price
Index in USDS-M Futures'', pages d'aide~; point d'accès public du taux de financement. Consulté le
septembre 2026.
\item BitMEX, point d'accès à la description des instruments de l'API publique (XBTUSD, ETHUSD, XBTUSDT), 24 septembre 2026.
\item M.~Schmeling, A.~Schrimpf et K.~Todorov, ``Crypto carry'', BIS Working Papers 1087, avril 2023.
\item D.~Ackerer, J.~Hugonnier et U.~Jermann, ``Perpetual Futures Pricing'', document de travail, 27 février 2025.
\end{omsources}

\section{Exercices}

\begin{exercise}[$\star$]\label{exo:m3:perpetual-futures:1}
Une position acheteuse de 2 bitcoins dans un \omterm{def:m3:perpetual-futures:perp}{contrat perpétuel} linéaire est détenue lors d'un moment de financement, le \omterm{def:m3:perpetual-futures:index}{prix de marquage} étant de 60\,000 et le \omterm{def:m3:perpetual-futures:perp}{taux de
financement} de 0.03~\%. Qui paie qui, et combien~?
\end{exercise}

\begin{exercise}[$\star$]\label{exo:m3:perpetual-futures:2}
La composante d'intérêt étant de 0.01~\% et le plafonnement de 0.05~\%, quel est le \omterm{def:m3:perpetual-futures:perp}{taux de financement} pour des primes moyennes de 0.02~\%,
0.10~\% et $-0.08\%$~?
\end{exercise}

\begin{exercise}[$\star$]\label{exo:m3:perpetual-futures:3}
Pourquoi la plateforme utilise-t-elle une médiane de plusieurs prix pour le \omterm{def:m3:perpetual-futures:index}{prix de marquage} plutôt que le dernier cours négocié~?
\end{exercise}

\begin{exercise}[$\star\star$]\label{exo:m3:perpetual-futures:4}
Un trader est acheteur de 60\,000 \omterm{def:m3:perpetual-futures:types}{contrats inverses} à 60\,000. Quel est le P\&L en bitcoins si le prix passe à 50\,000~? À
75\,000~? En dollars~?
\end{exercise}

\begin{exercise}[$\star\star$]\label{exo:m3:perpetual-futures:5}
Le \omterm{def:m3:perpetual-futures:types}{contrat quanto} ETHUSD verse 0.000001 bitcoin par contrat et par dollar du prix de l'ether. Un trader est acheteur de 10\,000 contrats
de 3\,000 à 3\,300~; le bitcoin vaut 60\,000. Quel est le P\&L en bitcoins et en dollars~?
\end{exercise}

\begin{exercise}[$\star\star$]\label{exo:m3:perpetual-futures:6}
Le financement reste à 0.01~\% toutes les huit heures pendant un an. Que perçoit un vendeur sur USD~10 millions~?
\end{exercise}

\begin{exercise}[$\star\star\star$]\label{exo:m3:perpetual-futures:7}
\emph{Codage.} À l'aide de \texttt{funding\_trade} et \texttt{yearly\_funding}, calculer, pour chaque année, le rendement sur
capital de la stratégie de financement. Pour quelles années est-il inférieur à un taux monétaire de 5~\%~?
\end{exercise}

\begin{exercise}[$\star\star\star$]\label{exo:m3:perpetual-futures:8}
\emph{Repérer l'erreur.} «~L'achat au comptant combiné à la vente du perpétuel est neutre au marché, donc il ne peut pas faire perdre d'argent.~»
\end{exercise}

\section{Problème~: la stratégie de financement}

\begin{problem}[{Problème du week-end --- encaisser un an de financement}]\label{pb:m3:perpetual-futures:1}
Un fonds achète pour USD~100 millions de bitcoins au comptant et vend le même notionnel d'un \omterm{def:m3:perpetual-futures:perp}{contrat perpétuel} linéaire, en payant des frais de preneur de
0.10~\% sur le comptant et de 0.05~\% sur le perpétuel, à l'entrée et à la sortie. Il conserve la position pendant un an au \omterm{def:m3:perpetual-futures:perp}{taux de financement} de 2024
(\cref{dat:m3:perpetual-futures:years}), le taux de marge de maintenance étant de 0.5~\%.

\medskip\noindent\textbf{Partie I --- Le revenu.}
\begin{enumerate}
\item Quel financement la vente perçoit-elle sur l'année, en supposant le prix constant~?
\item Quels sont les frais d'entrée et de sortie~?
\item À combien s'élève sur un an le financement par défaut de 0.01~\% par intervalle~?
\item Pourquoi le financement moyen de 2024 était-il proche de ce défaut~?
\item Quel aurait été le revenu au financement de 2022~?
\end{enumerate}

\medskip\noindent\textbf{Partie II --- La marge.}
\begin{enumerate}[resume]
\item Quelle marge la jambe vendeuse doit-elle détenir pour survivre à une hausse de 30~\%~?
\item Quel est le capital de la stratégie, et quel est son rendement sur capital~?
\item Quel est le rendement sur capital si la plateforme compte le comptant comme collatéral de la vente, et qu'aucune marge
supplémentaire n'est requise~?
\item Que se passe-t-il si le prix monte de 40~\% et que la marge n'est pas réapprovisionnée~?
\item Pourquoi le gain du comptant pendant la hausse ne compense-t-il rien, à moins que la plateforme ne compense les deux positions en marge croisée~?
\end{enumerate}

\medskip\noindent\textbf{Partie III --- Les risques.}
\begin{enumerate}[resume]
\item Que devient la stratégie lorsque le financement devient négatif~?
\item Quelle est l'exposition de la stratégie à la plateforme~?
\item Quelle est son exposition au \omterm{def:m3:blockchains-for-traders:stable}{stablecoin} dans lequel le perpétuel est margé~?
\item Comment la base à la sortie influe-t-elle sur le résultat~?
\item Que signifie pour cette stratégie le résultat de la BIS selon lequel un carry élevé annonce des krachs~?
\end{enumerate}

\medskip\noindent\textbf{Partie IV --- Jugement.}
\begin{enumerate}[resume]
\item La stratégie de financement est-elle un arbitrage~?
\item Pourquoi le financement existe-t-il, du point de vue de ceux qui le paient~?
\item Comment on dimensionnerait la stratégie sur plusieurs plateformes~?
\item Énoncer le \emph{résultat nommé}~: le rendement sur capital de la stratégie de financement sur 2024 avec 30.65~\% de marge, et la
marge nécessaire pour survivre à une hausse de 30~\%.
\item En une phrase~: que procure un paiement de financement à l'acheteur~?
\end{enumerate}
\end{problem}

\section{Questions d'entretien}

\begin{interviewq}[$\star$ \iqroles{trader}]\label{iq:m3:perpetual-futures:1}
Qu'est-ce qu'un \omterm{def:m3:perpetual-futures:perp}{contrat perpétuel} et qu'est-ce qui le maintient proche du comptant~?
\end{interviewq}

\begin{interviewq}[$\star$ \iqroles{développeur}]\label{iq:m3:perpetual-futures:2}
Pourquoi les plateformes liquident-elles sur le \omterm{def:m3:perpetual-futures:index}{prix de marquage} plutôt que sur le dernier cours négocié, et comment ce prix est-il construit~?
\end{interviewq}

\begin{interviewq}[$\star\star$ \iqroles{chercheur, trader}]\label{iq:m3:perpetual-futures:3}
Expliquez la différence entre les \omterm{def:m3:perpetual-futures:perp}{contrats perpétuels} linéaires, inverses et quanto, et dites qui utiliserait chacun.
\end{interviewq}

\begin{interviewq}[$\star\star$ \iqroles{trader}]\label{iq:m3:perpetual-futures:4}
Le financement est de 0.01~\% depuis des semaines alors que le perpétuel se traite à quelques points de base au-dessus de l'indice. Pourquoi~?
\end{interviewq}

\begin{interviewq}[$\star\star$ \iqroles{risque}]\label{iq:m3:perpetual-futures:5}
Que peut-il mal tourner avec un livre acheteur au comptant et vendeur sur perpétuel de USD~500 millions~?
\end{interviewq}

\begin{interviewq}[$\star\star\star$ \iqroles{développeur, chercheur}]\label{iq:m3:perpetual-futures:6}
Vous devez prédire le prochain \omterm{def:m3:perpetual-futures:perp}{taux de financement} deux heures avant le moment de financement. Comment procéderiez-vous~?
\end{interviewq}
